Kapitola shrnuje výsledky simulační studie. Zaměřujeme se na dvě klíčové metriky:
\begin{enumerate}
    \item \textbf{Relative Bias (Relativní zkreslení)}: O kolik procent Woodruffův odhad rozptylu nadhodnocuje nebo podhodnocuje skutečný rozptyl výběrového kvantilu (zjištěný empiricky z 1000 replikací).
    \item \textbf{Coverage Probability (Pokrytí)}: Podíl intervalů spolehlivosti, které obsahovaly skutečnou hodnotu $Q_p$. Očekávaná hodnota je 0.95.
\end{enumerate}

\section{Analýza přesnosti Woodruffova odhadu}

Tabulka 4.1 ukazuje výsledky pro $\sigma=0.5$ (velmi mírná asymetrie). Vidíme, že metoda funguje excelentně. Pro $n=500$ je pokrytí pro medián 0.946 a pro 90. percentil 0.959, což je blízko nominální hladině.

\begin{table}[h]
\centering
\caption{Výsledky simulace pro Log-Normal(0, 0.5) a Log-Normal(0, 1.5)}
\begin{tabular}{|l|l|l|r|r|}
\hline
\textbf{Sigma} & \textbf{n} & \textbf{Kvantil} & \textbf{Bias (\%)} & \textbf{Pokrytí} \\ \hline
0.5 & 200 & Median & 8.5 \% & 0.946 \\
0.5 & 200 & P90 & 32.6 \% & 0.953 \\
0.5 & 2000 & Median & 10.0 \% & 0.955 \\
0.5 & 2000 & P90 & 14.0 \% & 0.952 \\ \hline
1.5 & 200 & Median & 5.0 \% & 0.938 \\
1.5 & 200 & P90 & \textbf{77.5 \%} & 0.959 \\
1.5 & 2000 & Median & 3.5 \% & 0.946 \\
1.5 & 2000 & P90 & 6.3 \% & 0.946 \\ \hline
\end{tabular}
\end{table}

Při zvýšení šikmosti na $\sigma=1.5$ (spodní část tabulky) sledujeme zajímavý jev.
U 90. percentilu pro $n=200$ je relativní bias obrovský (77.5 \%), což značí, že Woodruffova metoda silně nadhodnocuje rozptyl. Přesto je pokrytí (Coverage) konzervativní (0.959). To je způsobeno asymetrií skutečného výběrového rozdělení, kterou metoda nahrazuje symetrickým intervalem (nebo v našem případě je interval odvozen z robustních kvantilů, což ho činí širokým).

Se zvyšujícím se rozsahem výběru ($n=2000$) bias klesá k přijatelným hodnotám (6.3 \%) a pokrytí konverguje k 0.95.

\section{Zhodnocení vlivu rozsahu výběru}

Simulace potvrdila teoretický předpoklad, že pro odhady v "řídkých" oblastech hustoty (chvosty log-normálního rozdělení) konverguje rozptyl k asymptotickému tvaru pomaleji. 
Zatímco pro medián stačí i $n=200$ k velmi spolehlivým výsledkům, pro 90. percentil silně sešikmeného rozdělení dává Woodruffova metoda při malých rozsazích výběru spíše konzervativní (příliš široké) intervaly, což vede k nadhodnocení rozptylu, ale "bezpečnému" pokrytí.
